% ============================================================
% main_v2.tex  短縮・再構成版の親ファイル（2026-10-27）
%   方針：論文_再構成案.md（改訂版）。投稿先の第一候補は Adaptive Behavior（日本語で固めてから英訳）。
%   本文は sections_v2/*.tex，補足資料は supplement_v2.tex（supp_v2/*.tex）。
%   図は figs/，表は tables/，文献は references.bib を初稿と共有する。
% コンパイル：upLaTeX（latexmkrc は初稿と同じ）。Overleaf で main_v2.tex を主文書にする。
% ============================================================
\documentclass[uplatex,dvipdfmx,a4paper,11pt]{jsarticle}

\usepackage{amsmath,amssymb}
\usepackage{graphicx}
\usepackage{booktabs}
\usepackage{array}
\usepackage{multirow}
\usepackage{xcolor}
\usepackage{url}
\usepackage[hidelinks]{hyperref}
\usepackage{pxjahyper}

% 図のファイルがないときは枠と名前だけを出す（図がそろったら消してよい）。
\makeatletter
\let\orig@includegraphics\includegraphics
\renewcommand{\includegraphics}[2][]{%
  \IfFileExists{#2.pdf}{\orig@includegraphics[#1]{#2}}{%
  \IfFileExists{#2.png}{\orig@includegraphics[#1]{#2}}{%
  \IfFileExists{#2}{\orig@includegraphics[#1]{#2}}{%
    \fbox{\parbox[c][4em][c]{0.9\linewidth}{\centering\ttfamily 図未作成：\detokenize{#2}}}}}}}
\makeatother
% 図表の置き場所の制約をゆるめる（大きな図が [t] に入らず，以降の図表がすべて最後に回るのを防ぐ）。
\renewcommand{\topfraction}{0.9}
\renewcommand{\bottomfraction}{0.8}
\renewcommand{\textfraction}{0.1}
\renewcommand{\floatpagefraction}{0.75}
\setcounter{topnumber}{3}
\setcounter{bottomnumber}{2}
\setcounter{totalnumber}{5}

\newcommand{\todo}[1]{\textcolor{red}{#1}}

\title{自分の価値の経験から学んだ予測は，どのような条件で他者の予測に転用されるか\\
\large ──自他で重みを共有する予測学習における成立条件の切り分け──}
% 英題（案）：When does prediction learned from one's own value-guided experience
%            transfer to predicting others? Dissecting the conditions in a self--other
%            shared predictive-learning model
\author{草野 \todo{【要確認：氏名・所属・共著者】}}
\date{}

\begin{document}
\maketitle

\begin{abstract}
シミュレーション説は，人が自分の心の働きを用いて他者を理解すると考える．
本研究は，自他で重みを共有する予測学習の仕組みである重ね合わせ機構の上で，自分の価値の経験から学んだ予測傾向が，どのような条件で他者側の予測に転用されるかを調べた．
方向ごとの快・不快を受けて次の視界を予測するエージェントを，価値に従って動く経験の量を変えて学習させた．
価値に従う行動の見込みは報酬を用いない予測学習で生まれ，他者側に与えた快・不快にも他者の予測が応じるようになった．
同じ軌跡で快・不快の符号だけを無作為にすると，自分側の快・不快への依存と他者側の符号に従う応答は0前後になり，他者側の符号応答も大きく弱まった．
快・不快の入口の重みを写すだけでは転用は戻らなかった．
一方，転用は他者の行動の予測を一般に正確にはせず，他者の行動の仕方が自分と違うと本当の快・不快を与えても予測が外れることがあり，他者固有の好みも書き出されなかった．
成立条件を切り分けた対照は，自己位置の追跡の事前基準を満たさなかった条件で行った．
\end{abstract}

\noindent\textbf{キーワード}：シミュレーション説，予測学習，自他の共有表現，価値，計算モデル

% ============================================================
% v2/sections/intro.tex  はじめに（短縮・再構成版 2026-10-27）
% 方針：論文_再構成案.md（改訂版）。主の貢献は「成立条件の切り分け」。A は補助1（可能性）。
% ============================================================
\section{はじめに}\label{sec:intro}

シミュレーション説は，他者の立場に自分を置き，自分の心の働きを用いて他者の心的状態や行動を推し量ると考える\cite{gordon1986folk,goldman2006simulating}．
予測符号化の観点からは，自分の意思決定や計画の仕組みが，他者の行動を説明する仮説を試すためにも使われ，そこでの信念や欲求は自他を区別しない形で表されると論じられている\cite{gordon2021simulation}．
しかし，自分の価値に基づく予測の仕組みが，どのような条件で他者の予測に使われるようになるのかを，学習の過程から具体的に示した計算モデルは少ない．

Noguchi ら\cite{noguchi2022superposition}の重ね合わせ機構（superposition mechanism）は，一つの感覚入力を二つの経路で符号化し，重みを共有した一つのモジュールで処理する予測学習の仕組みである．
その新しさは，他者の動きを自分のどの経験に当てはめるかという自他の対応を与えずに，自分の次の視界を予測する学習だけから，自他の区別，他者の位置，他者の視点，他者の運動が生まれた点にある．
一方で，この機構が扱ったのは他者の位置・視点・運動までであり，他者の欲求のような内面には及んでいない\cite{noguchi2022superposition}．

本研究は，重ね合わせ機構の上で，自分の価値に従って動いた経験から学ばれた予測が他者の予測に使われるかを調べる．
ここで価値は，方向ごとの快・不快として人が与える信号である（\ref{sec:env}節）．
問うのは価値そのものが生まれるかではなく，与えた価値の信号とそれに従う行動の規則性を，自他で共有された予測学習がどう利用するかである．
また本研究は，自他の対応を担う視覚エンコーダを，文献\cite{noguchi2022superposition}の設定で学習した重みに固定して与える．
自他を同じ符号で表す性質は本研究の新しさではなく，前提である．

重みを共有している以上，自分側で学んだことが他者側にも働くのは当然だ，と考えられるかもしれない．
しかし，他者側に届くのが「入力された快・不快の信号」そのものなのか，「その信号が自分の動きの何を予測するか」として学ばれた解釈なのか，その解釈が仕組みのどこに宿るのかは自明でない．
そこで本研究は，次の問いを立てる．
\begin{quote}
  自分の価値の経験から学んだ予測は，自他で共有された予測の仕組みを通して，どのような条件のもとで他者の予測に使われるか．
\end{quote}
ここで「使われる」（転用）とは，他者側に与えた価値の入力に応じて他者の予測が変わることを指し，他者の行動の予測が正確になることとは区別する．
これを，価値に従って動いた経験の量を変えた学習と，三つの対照で調べた．
第一に，自他で重みを共有しない対照（非共有条件）．
第二に，自分の動きの軌跡はまったく同じまま，快・不快の符号だけをエピソードごとに無作為にし，価値の信号と行動の対応だけを壊した対照（符号無作為条件）．
第三に，学習を終えた非共有のモデルに，自分側の重みを写す比較である．

本研究の貢献は次のとおりである．
\begin{description}
  \item[主] 価値の信号と行動の対応から学習された自分の予測傾向が，他者側の予測に反映されること（転用）を示し，その成立条件を切り分けた．同じ軌跡で価値の信号と行動の対応を壊すと，自分側の快・不快への依存と他者側の符号に従う応答は0前後になり，他者側の符号応答も大きく弱まった．他者の状態を追う力が育つだけでは転用は生じず，学習後に快・不快の入口の重みだけを写しても戻らなかった．転用には，価値の信号と行動の一貫した対応と，自分の経験で学習された入口以外の重みや計算が要ると考えられる．
  \item[補助1] この転用は，他者の行動の予測を一般に正確にするものではない．本当の快・不快を与えたときの誤差は経験によって小さくならず，他者の行動の仕方が自分と違うと，経験の少ないモデルより外れることがあった．他者の規則の確率を自分の学習条件に近づけると，本当の快・不快の役立ち方が変わったが，位置の分布も同時に変わるので，自分の行動の仕方の投影という解釈は可能性にとどまる．
  \item[補助2] この転用は，他者固有の好みの自発的な表現や推定を保証しない．他者用のモジュールは好みではなく位置や方角を書き出し，好みの候補を入れ比べる照合も，他者の進んだ方角だけで選ぶ基準をほとんど上回らなかった．
\end{description}

なお，成立条件を切り分けた対照は，自己位置の追跡の事前基準を満たさなかった経験80\%条件で行った．主張の強さへの影響は\ref{sec:disc-strength}節で論じる．

以下，\ref{sec:related}章で関連研究を整理し，\ref{sec:method}章でモデルと実験方法を，\ref{sec:results}章で結果を述べる．
\ref{sec:discussion}章でシミュレーション説との関係と限界を論じ，\ref{sec:conclusion}章でまとめる．

% ============================================================
% v2/sections/related.tex  関連研究（短縮・再構成版 2026-10-27）
% ============================================================
\section{関連研究}\label{sec:related}

\paragraph{シミュレーション説と自他の共有表現}
神経科学では，自分の行為・感覚と他者のそれとに共通する表現と，両者を区別する仕組みの両方が社会的認知に必要だと整理されている\cite{decety2003shared}．
計算論の側では，運動制御のための内部モデルを他者の行為の理解に用いる枠組み\cite{wolpert2003unifying}や，ミラーニューロン系の働きを予測符号化として捉える考え方\cite{kilner2007predictive}が提案されてきた．
Gordon\cite{gordon2021simulation}は，自分の意思決定の仕組みが他者の行動の仮説を試すために使われ，そこでの信念や欲求は自他を区別しない符号化で表されると論じた．
重ね合わせ機構\cite{noguchi2022superposition}は，こうした考えを，自他の区別を与えない予測学習の計算モデルとして具体化したものと位置づけられる．
本研究は，この機構の上で，自分の価値に関わる予測が他者に使われるための条件を調べる．

\paragraph{投影と調整}
人は，自分の判断や好み，その時の身体の状態を他者に当てはめて過大に見積もる\cite{ross1977false,nickerson1999how,vanboven2003social}．
他者の視点を取ることは，自分の視点を出発点にして調整する過程として説明され\cite{epley2004perspective}，他者の好みを推し量るときの内側前頭前野の活動も，自分と他者の隔たりに応じて変わる\cite{tamir2010neural}．
Suzuki ら\cite{suzuki2012learning}は，他者の価値に基づく選択を予測するとき，自分の価値づけの信号が流用されること，他者の行動の予測誤差による補正が他者の価値が自分と異なるときに重要になることを示唆した．
本研究のモデルは，この区別でいえば自分の価値づけを流用する経路を持ち，他者の価値を補正する専用の経路を持たない．

\paragraph{自分をもとに他者を予測する計算モデル}
Raileanu ら\cite{raileanu2018modeling}は，自分の目標を入力の一つとする自分の方策を他者に当てはめ，他者の行動がもっともらしくなるように目標の入力を最適化して他者の目標を推した．
そこでは自分のモデルは報酬に基づく強化学習で得た方策で，目標が入力として明示されている．
本研究の見込みは，報酬も目標の入力も用いない予測学習の中で，自分の行動の経験から生まれる．
Copete ら\cite{copete2016motor}は，一つの予測器を自分の行為と他者の行為の観察の両方に用い，自分の運動の学習が他者の行為の結果の予測を助けることを示したが，観察者と行為者が同じ視点を共有すると仮定して自他の対応を与えている．
重ね合わせ機構では他者の内部状態が自分とは別に保たれるので，自他の重みの共有だけを外す対照や，学習後に自分側の重みを写す比較を作ることができる．
本研究はこれを用いて，他者の予測に使われるのが何かを切り分ける．
他者の行動から目標や好みを推す問題は，逆計画\cite{baker2009action,baker2017rational}や逆強化学習\cite{ng2000algorithms}として広く研究されている．
本研究の照合も候補ごとの予測を観察と比べる構造を持つが，推論の手法そのものを新しいとは主張しない．

% ============================================================
% v2/sections/methods.tex  モデルと実験方法（短縮・再構成版 2026-10-27）
% 数字は sections/methods.tex（第2弾の照合済み）と Claude Code の報告から。
% ============================================================
\section{モデルと実験方法}\label{sec:method}

\begin{figure}[tbp]
  \centering
  \includegraphics[width=0.49\linewidth]{figs/env_conditions}\hfill
  \includegraphics[width=0.49\linewidth]{figs/network}
  % 左：環境と条件（env_conditions.pdf）。右：ネットワークと2段階の学習（network.pdf）。
  % 投稿の前に一枚の図に作り直すことを検討（Claude Code）。
  \caption{環境とモデル．左：アリーナと赤・緑を入れ替えた配置，全方位の視界と方向ごとの快・不快 $P_k$，段階Iの学習データにおける A-1 の滞在位置．右：ネットワーク．視覚エンコーダ1・2（斜線，凍結）が A-1 の視界を読み，重みを共有した LSTM の process-1（自分）と process-2（他者）が別々の内部状態 $h^1$，$h^2$ を持つ．差し込みは非共有条件で分ける部分．}
  \label{fig:model}
\end{figure}

\subsection{環境と価値の信号}\label{sec:env}
環境は文献\cite{noguchi2022superposition}と同じ $20\times20$ のアリーナで，四隅の外に赤・緑・青・シアンの立方体を置いた（図\ref{fig:model}）．
エージェントは2体で，学習するのは A-1 だけである．
A-1 は回転しない全方位カメラで $16\times64$ 画素の視界 $v_t$ を得る．
文献\cite{noguchi2022superposition}との違いとして，エピソードごとに確率 1/2 で赤と緑の位置を入れ替えた．
入れ替えのない環境では，「赤の方へ動く」ことと「特定の場所へ動く」ことが区別できないためである．
各エピソードは 101 フレームで，A-1 は1ステップに1の距離を動くか，休止（確率 0.005 で始まり 10〜30 ステップ）する．

A-1 の価値は，視界から毎ステップ計算する8次元の快・不快 $P_t$ として与えた．
カメラの8列ごとの範囲 $k$ で $P_{t,k}=(n^{\mathrm{like}}_{t,k}-n^{\mathrm{dislike}}_{t,k})/128$ とし，好きな色を赤，嫌いな色を緑とした（$n$ は各色の画素数）．
値が0以外になるのは範囲 0・1・6・7 だけである．
この値は人が与える信号であり，ネットワークが学ぶのは，この信号と自分の行動・視界の変化との関係である．

価値に従う一歩は，視界の中の赤の円周上の重心の方角へ進む規則で作った（壁では沿って進み，進めないときはランダム）．
A-1 は休止していない各ステップで，確率 $p$ で規則の一歩を，残りでランダムな一歩をとる．
$p=0$ を経験なし条件，$p=0.5$ を経験50\%条件，$p=0.8$ を経験80\%条件と呼ぶ．
$p$ は規則の一歩を選ぶ確率で，実際に規則で動いた一歩の割合は経験50\%条件で 48\%，経験80\%条件で 57\% であった．
乱数は，同じエピソード番号で配置・初期位置・休止の時刻が条件の間で一致するように設定した．
ただし $p$ を上げると位置の分布も変わり，赤から距離5以内にいるフレームの割合は約 0.03，約 0.33，約 0.85 であった．
この交絡を扱うため，後述の符号無作為条件を作った．

段階Iでは A-2 は静止し，各条件で学習用 3000・評価用 300 エピソードを集めた．
段階IIの学習と，他者側の評価のうち符号に従う分・真値条件・照合では，A-1 はランダムに歩き，A-2 は自分の視界で同じ規則を用いて一歩ごとに確率 0.5 で好きな色へ向かう．
A-2 が常に緑を好むデータ（緑好みデータ）と，エピソードごとに緑好き・赤好きが変わるデータ（好み混合データ）を用いた．

\subsection{ネットワーク}\label{sec:network}
ネットワークは文献\cite{noguchi2022superposition}に従う（図\ref{fig:model}）．
二つの視覚エンコーダが A-1 の視界を読み，共有モジュール（128 ユニットの LSTM）が process-1 と process-2 で重みを共有して内部状態 $h^1_t$，$h^2_t$ を別々に持つ．
視覚予測器が両者をまとめて A-1 自身の次の視界 $v_{t+1}$ を予測し，特徴予測器が $h^1$，$h^2$ から次の視覚エンコーダの出力を予測する．
学習中は視覚を各ステップ確率 0.99 で隠す．
文献\cite{noguchi2022superposition}では process-1 の入力は運動 $m_t$ だったが，本研究では $x^1_t=[m_t,P_t,b_t]$（11次元）とした．
運動は 5〜15 ステップの区間を交互に置いて隠し（$t\ge1$ の約48\%），隠れた区間では $m_t=0$，隠れフラグ $b_t=1$ とし，$P_t$ を区間の最初の値に保つ．
運動が分からないときに次の視界を当てるには自分の動きを推し量る必要があり，価値に従って動いた経験があれば「快の高い方へ動いた」と推し量る理由が生まれる．
段階Iでは process-2 の入力は零である．

視覚エンコーダは，文献\cite{noguchi2022superposition}の設定（運動を入力，入れ替えなし）で再現学習した重み（基準エンコーダ）に，すべての条件で固定した．
基準エンコーダを用いた再現モデルでは，$h^1$ から A-1 の位置が $R^2=0.967$，$h^2$ から A-2 の位置が 0.993 で読め，視覚エンコーダ2と1の符号のコサイン類似度（自他一致度）は 0.80 である．
運動以外の入力を加えてエンコーダごと学習すると自他一致度が 0.18〜0.27 にとどまったため，固定した（補足資料）．
したがって自他を同じ符号で表す性質は学んだものではなく与えたものである．

段階Iは，200 エポックの二乗誤差（MSE 期間）で共有モジュール・視覚予測器・特徴予測器を学習し，続く 200 エポックの絶対誤差（L1 期間）で視覚予測器の復号部だけを学習した（AdamW，学習率 $10^{-3}$，重み減衰 0.01，バッチ 10）．
段階IIでは，文献\cite{noguchi2022superposition}の Motion Generator と同じ構成の他者用モジュール $G$ を加え，マスク前の視覚エンコーダ2の出力から8次元の $Q2_t$ を作って process-2 の快・不快の枠に入れ（$x^2_t=[0,Q2_t,1]$），$G$ 以外を凍結して予測誤差で学習した．

\subsection{条件と対照}\label{sec:conditions}
条件と組の一覧，各指標の位置づけを表\ref{tab:design}に示す．
同じシードで学習した各条件のモデルを一つの「組」とし，同じ組の中では初期値・バッチの順・運動を隠す区間を一致させた．
事前登録した判定は本番の3組（組1〜3）で行い，予備の組4と追加の組5・6 も除外せずに報告する．

\paragraph{非共有条件}共有モジュールの LSTM を process-1 用と process-2 用に分け，経験80\%条件のデータで学習した（組1〜3）．process-2 の快・不快を受ける入力の重みは一度も学習されないので，他者側の働きが出ないことは構成からほぼ予想される．

\paragraph{符号無作為条件}経験80\%条件と全く同じ軌跡のデータで，$P_t$ の符号だけをエピソードごとに確率 1/2 で反転した（組1〜3．反転は学習用 3000 のうち 1471 エピソードで，配置とは独立）．位置・運動・視界の系列は経験80\%条件と同じで，価値の信号と行動の対応だけが壊れる．ただし赤へ近づくという行動の履歴そのものは変わらない．

\paragraph{写しの比較}学習を終えた非共有条件のモデルの process-2 用の LSTM に，process-1 用の LSTM の重みを写した．すべて写す場合と，入力の重みのうち快・不快を受ける8列だけを写す場合を作り，経験なし条件の非共有モデルにも同じように写した（組1〜3）．

\paragraph{他者の規則の確率}段階IIの評価用データと同じ作りで，A-2 が確率 0.8 で好きな色へ向かうデータを作った（A-1 の歩き・休止・配置・好きな色は元と全エピソードで一致）．

\begin{table}[tbp]
  \centering
  \footnotesize
  \caption{条件，指標，事前登録の区別}
  \label{tab:design}
  % 仮のファイル：Claude Code が本物の表で上書きする（tables/conditions.tex と metrics.tex を一枚にまとめる）。
\fbox{\parbox{0.9\linewidth}{\centering\ttfamily 表未作成：tables/v2\_design}}

  % 表1（新しく作る，Claude Code）：上段に条件（経験なし・50%・80%，非共有，符号無作為，写し）と組の数，
  % 下段に指標（G3，d_P，G5b，符号に従う分，真値条件の Δ，照合の BA，T2・T5）と位置づけ
  % （事前登録した主／事後に定義／事前登録外）。今の tables/conditions.tex と metrics.tex を一枚にまとめる。
\end{table}

\subsection{測る量}\label{sec:evaluation}
\paragraph{自分側の見込み}運動と視覚を $t=20$ から15ステップ隠し，快・不快を $t=20$ の値に保つ．各ステップの予測を「$t=20$ の位置から色 $X$ の方へ $d$ 歩動いた視界」と比べて最もよく合う $d$ を求め，その傾き $s_X$ から $g=s_{\text{赤}}-\frac13(s_{\text{緑}}+s_{\text{青}}+s_{\text{シアン}})$ を求めた（G3）．運動が分からないときに予測が赤の方へだけ進み続ければ $g>0$ になる．快・不快の符号を反転して保った場合の $g$ も求め，$d_P=(g_{\text{保つ}}-g_{\text{反転}})/2$ を自分側の見込みの快・不快への依存とした．

\paragraph{他者側への働き}主の指標（事前登録）は G5b である．A-2 が静止した評価で，$t=20$〜29 に process-2 の一つの範囲 $k$ だけへ $\pm a$（$a=0.28$）を入れ，10ステップ後の予測画像の中の A-2 の位置のずれの差 $\Delta_k(+a)-\Delta_k(-a)$ を範囲 $k$ の方向に射影して8範囲で平均した．同じ入力の履歴からの差なので，A-2 をどこに置いていたかに左右されない．
二つ目は，A-2 の本当の快・不快を保って入れた場合と符号を反転して入れた場合の，予測の中の A-2 の好きな色の方への傾きの差の半分（符号に従う分）である．この指標は，事前登録した関門 G5a に戻る動きが混ざることが分かった後に定義した．組1〜4 では事後に定義した指標で，事前登録してから計算したのは組5・6 だけである．

G5b と符号に従う分は，他者側の予測が価値の入力に応じて変わるか（転用）を測る量であり，他者の行動の予測が正確になるかは測らない．

\paragraph{本当の快・不快の役立ち方}予測の正確さへの効果は次のように測った．process-2 に A-2 の本当の快・不快を入れた場合（真値条件）と零を入れた場合で，他者側の特徴予測の誤差を比べ，差 $\Delta$（負なら本当の快・不快が役立つ）を求めた．この量は特徴予測器を通り，視覚予測器を通らない．段階Iの process-2 には零しか入らないので，本当の快・不快は学習で見ていない入力である．

\paragraph{他者の好みの表現と推定}他者用モジュールについては，赤と緑を途中で入れ替えたときに $Q2$ が新しい緑の範囲へ移るか（T2）と，位置をそろえた緑好き・赤好きの A-2 で $Q2$ が好みで違うか（T5）を，事前登録した見分けとした．
照合では，好み混合データで process-2 に二つの好みの候補（範囲 6・7 が快い形 $c^+$ とその逆 $c^-$）を全ステップ入れ，予測誤差の小さい方を選んだ．主の誤差は A-1 の視界の予測誤差（自己視界基準，L1 期間），副は他者側の特徴予測の誤差（他者特徴基準，MSE 期間）とし，balanced accuracy（BA）を求めた．
A-2 の好みはその間に進んだ方角とも結びつくので，A-2 の $t=1$〜50 の変位の向きだけで選ぶ基準（方角だけの基準）を並べた．

\paragraph{判定}各モデルを学習の最後の5つのチェックポイントで評価し，同じエピソードを対にしたブートストラップ（1000回）で，経験ありの値と組の差（経験あり $-$ 対照）の 95\% 信頼区間を求めた．
どちらも仮説の向きに0を含まない点が5点中3点以上ならその組を合格とし（照合の BA では経験ありの値の下限が 0.5 を上回ること），3組中2組以上で合格とした．
最終5点は互いに相関するので独立な繰り返しではなく，不確かさの主な単位は組である．
そのため組をまたいだ差と $t$ の 95\% 信頼区間（事前登録外の記述）を並べる．
判定の基準は計算の前に記録した（事前登録）．事後に定義した指標と事前登録外の確認は，そのように明記する．多重比較の補正は行っていない．

% ============================================================
% v2/sections/results.tex  結果（短縮・再構成版 2026-10-27）
% 数字は sections/results.tex（第2弾の照合済み）と Claude Code の報告（A・B・C）から。
% 組1=s0，組2=s1，組3=s2（本番），組4=s3，組5=s4，組6=s5。
% ============================================================
\section{結果}\label{sec:results}

\begin{figure}[tbp]
  \centering
  \includegraphics[width=\linewidth]{figs/central}
  % 元：paper/figs/central.pdf（G3，符号に従う分（緑好み・好み混合），G5b。6組，非共有は白い菱形）
  \caption{自分側の見込みと他者側への転用．横軸は経験の量．細い線は各組（組1〜3 は灰の実線，組4 は橙の破線，組5・6 は橙の点線と白抜きの四角），太い線は6組の平均と95\%信頼区間，白い菱形は非共有条件（組1〜3）．(a) 自分側の積み上がり $g$（事前登録）．(b)(c) 他者側の符号に従う分（緑好みデータ，好み混合データ．組1〜4 では事後に定義）．(d) 他者側の符号応答 G5b（事前登録）．}
  \label{fig:central}
\end{figure}

\begin{table}[tbp]
  \centering
  \footnotesize
  \setlength{\tabcolsep}{4pt}
  \caption{組をまたいだ主な差（全6組．各組の最終5点の平均を1つの値とし，組どうしを対にした差の平均と $t$ の95\%信頼区間．事前登録外の記述）}
  \label{tab:summary}
  \begin{tabular}{llll}
    \toprule
    指標 & 50\% $-$ 経験なし & 80\% $-$ 50\% & 80\% $-$ 経験なし \\
    \midrule
    自分側の積み上がり $g$ & $+0.071$ [$+0.064$, $+0.079$] & $+0.141$ [$+0.131$, $+0.152$] & $+0.213$ [$+0.205$, $+0.221$] \\
    G5b & $+0.143$ [$-0.109$, $+0.394$] & $+0.218$ [$-0.066$, $+0.503$] & $+0.361$ [$+0.057$, $+0.665$] \\
    符号に従う分（緑好み） & $+0.138$ [$+0.028$, $+0.248$] & $+0.207$ [$+0.050$, $+0.363$] & $+0.345$ [$+0.132$, $+0.557$] \\
    符号に従う分（好み混合） & $+0.137$ [$+0.027$, $+0.246$] & $+0.163$ [$+0.059$, $+0.267$] & $+0.300$ [$+0.174$, $+0.425$] \\
    \bottomrule
  \end{tabular}
  % 出典：v13 まとめ §11（第2弾の照合済み）。組ごとの値は補足資料の表。
\end{table}

\subsection{自分側の見込み}\label{sec:res-self}
自分側の積み上がり $g$ は経験の量に応じて段階的に大きくなった（図\ref{fig:central}(a)，表\ref{tab:summary}）．
本番の3組の値は，経験なし条件で $+0.019$〜$+0.033$，経験50\%条件で $+0.089$〜$+0.103$，経験80\%条件で $+0.225$〜$+0.240$ であった．
事前登録した組ごとの判定では，経験なし条件との差は経験50\%条件・経験80\%条件とも6組すべてで成り立ち，経験80\%条件が経験50\%条件を上回ることも6組すべてで成り立った．
経験なし条件の $g$ も小さく正であり，見込みは経験なし条件よりはっきり強く生まれた，と言うのが正確である．

前提とした自己位置の追跡（運動が見えるステップで $h^1$ から A-1 の位置を読む決定係数 0.80 以上）は，経験50\%条件では6組とも満たしたが，経験80\%条件では6組とも満たさなかった（0.70〜0.75）．
自分の学習データと同じ作りのテストでは 0.99 であり，赤の角にとどまる学習データによる，経験していない場面への一般化の弱さと考えられる．
$g$ などの指標は位置の読み出しを使わずモデル自身の予測から測るので計算はできるが，経験80\%条件の値はこの状態を含むものとして読む必要がある．
前提を満たした経験50\%条件と経験なし条件の比較だけでも，経験によって見込みが生まれることは6組で成り立つ．
非共有条件でも自分側の見込みは経験80\%条件と同じように生まれた（$g$ は $+0.238$〜$+0.248$）．

\subsection{他者側への転用}\label{sec:res-other}
process-2 の一つの範囲に $\pm a$ を入れると，経験80\%条件では予測の中の A-2 がその向きにずれた（図\ref{fig:central}(d)）．
本番の3組では，経験80\%条件の G5b は $+0.392$〜$+0.682$（経験なし条件 $-0.191$〜$-0.030$）で，事前登録した判定に3組とも合格した．
全6組では経験なし条件との差が5組で成り立ち，組をまたいだ差も0を含まなかった（表\ref{tab:summary}）．
自己位置の追跡の前提を満たした経験50\%条件でも，G5b は $+0.063$〜$+0.197$ で，事前登録した判定に合格した（本番の3組中2組）．
ただし組をまたいだ段階ごとの差は0を含み，G5b で組をまたいで示せるのは経験80\%条件が経験なし条件より大きいことまでである．
外れた組は，経験なし条件の G5b が大きかった組4（$+0.291$）と，経験50\%条件が経験80\%条件を上回った組6である（補足資料）．
A-2 の本当の快・不快を用いる符号に従う分では，二つのデータとも二つの段階の差が0を含まなかったが（表\ref{tab:summary}），この指標は組1〜4 では事後に定義したものである（\ref{sec:evaluation}節）．
これらは，自分側で学んだ傾向が他者側の予測に反映されること（転用）を示すが，他者の行動の予測が正確になることは示さない（\ref{sec:res-projection}節）．

非共有条件では他者側の働きは消えた（G5b は3組とも $-0.01$〜$+0.01$，符号に従う分も0前後）．
それでも $h^2$ から A-2 の位置を読む決定係数は 0.91〜0.93 で，他者の状態を追う力は育っていた．
他者の状態を追えることだけでは，他者側の価値に関わる働きは生じない．
ただしこの結果は，process-2 の快・不快の入口が学習されないという構成からほぼ予想されるものである．

\subsection{成立条件：価値の信号と行動の対応}\label{sec:res-signrand}

\begin{figure}[tbp]
  \centering
  \includegraphics[width=\linewidth]{figs/c_signrand}
  % 元：paper/figs/c_signrand.pdf。パネルの順を (a) d_P，(b) G5b，(c)(d) 符号に従う分，(e) g，(f) Δ（補助，薄く）に
  %   並べ替えてもらう（第3弾の指示書）。零の条件が加わったら差し替え。
  \caption{位置の分布をそろえた対照（符号無作為条件）．横軸は経験なし条件，符号無作為条件，経験80\%条件．主な量は (a) 自分側の快・不快への依存 $d_P$，(b) 他者側の符号応答 G5b，(c)(d) 符号に従う分（緑好みデータ，好み混合データ）．(e) 快・不快を保ったときの自分側の積み上がり $g$ と (f) 真値条件の $\Delta$ は補助．組1〜3 を線で結ぶ．}
  \label{fig:signrand}
\end{figure}

経験80\%条件と同じ軌跡で快・不快の符号だけを無作為にすると，自分側の快・不快への依存と他者側の符号に従う分は0前後になり，他者側の符号応答 G5b も大きく弱まった（図\ref{fig:signrand}）．
自分側の $d_P$ は0前後（経験80\%条件 0.23〜0.25），符号に従う分は二つのデータとも $-0.02$〜$+0.03$（経験80\%条件 $+0.18$〜$+0.44$）であった．
G5b は組1〜3 で $+0.15$，$+0.08$，$-0.05$ と，3組とも経験80\%条件（$+0.47$，$+0.68$，$+0.39$）より小さかったが，3組とも経験なし条件（$-0.05$，$-0.03$，$-0.19$）よりは大きく，完全に消えたとは言えない．
事前登録した予想のうち，他者側の働きが経験80\%条件より小さくなること（G5b・符号に従う分）と自分側の快・不快への依存が小さくなること（$d_P$）は，3組とも成り立った．
照合の他者特徴基準の BA も 0.50〜0.57 にとどまった（経験80\%条件 0.74〜0.78）．
一方，次の視界の予測誤差と自己位置の追跡は経験80\%条件に近く（追跡の決定係数は【要確認】），学習も安定していた（学習の損失は約5\%高いだけ）．
自分側の予測の質を保ったまま，快・不快に関わる働きが大きく弱まったことになる．
ただし，快・不快を保ったときの $g$ は経験なし条件より $+0.05$〜$+0.13$ 大きく残った．
これは事前登録で注意として挙げた，快・不快によらず赤へ向かう癖に当たる．

この対照は，位置・運動・視界の系列をそろえたまま，価値の信号と行動の一貫した対応だけを壊したものである．
したがって，快・不快に従う自分側の見込みと，それが他者側に及ぶ働きは，学習データの位置の偏りだけでは説明できず，価値の信号と行動の一貫した対応に大きく依存すると言える．
一方で，赤へ近づくという行動の履歴そのものは変えていないので，価値に従って動いた経験のあらゆる効果を識別したものではない．
G5b に残った分は，赤への癖や，位置の手がかりとしての快・不快の大きさを含みうる．
快・不快を常に零にした条件の結果は【結果待ち】である．

\subsection{成立条件：自分の予測の力学の再利用}\label{sec:res-copy}
学習を終えた非共有条件のモデルに自分側の重みを写すと，視覚予測器を通らない量では，共有ありの経験80\%条件と同じ形の働きが戻った（事前登録外の確認を含む．詳細は補足資料）．
LSTM の重みをすべて写すと，照合の他者特徴基準の BA は 0.75〜0.77 になり（元の非共有条件 0.41〜0.49，経験80\%条件 0.74〜0.78），経験なし条件の非共有モデルに同じように写した場合（0.35〜0.52）との差は3組とも正で信頼区間は0を含まなかった（$+0.25$〜$+0.41$）．
真値条件の $\Delta$ でも，経験ありの重みを写したときにだけ，経験80\%条件と同じく A-2 の規則の確率による変化が現れた．
快・不快を受ける入口の8列だけを写した場合は，G5b は0前後のままで（$-0.035$〜$+0.030$，事前登録した予想は不支持），照合も戻らなかった（0.45〜0.55）．
他者側への価値に関わる傾向の転用は，快・不快の入口の重みだけでは説明できず，自分の経験で学習された LSTM のほかの重みや計算にも依存する可能性が示された．
ただし，すべて写す操作は再帰の結合だけでなく入力の重みやバイアスも同時に変えるので，どの重みが要るかは特定していない．
すべて写すと他者の追跡が基準を下回り（$h^2$ の決定係数 0.60〜0.75），予測画像を通る量（G5b，符号に従う分）は，写した $h^2$ と視覚予測器の食い違いのために判定できなかった．
学習中に重みを共有することと学習後に写すことが同等であるとも言えない．

\subsection{転用は他者の予測を正確にするか}\label{sec:res-projection}

\begin{figure}[tbp]
  \centering
  \includegraphics[width=\linewidth]{figs/a_cross_delta}
  % 元：paper/figs/a_cross_delta.pdf
  \caption{他者の規則の確率を変えたときの，本当の快・不快の役立ち方．縦軸は真値条件の誤差から零を入れた誤差を引いた値 $\Delta$（負なら役立つ）．横軸は A-2 が規則の一歩を選ぶ確率 $p$（実際に規則で動いた一歩は $p=0.5$ で約46〜47\%，$p=0.8$ で約57\%）．(a) 緑好みデータ，(b) 好み混合データ．細い線は各組（組1〜3 は実線，組4〜6 は点線），太い線は全6組の平均．}
  \label{fig:cross}
\end{figure}

転用が他者の行動の予測を正確にするかを，A-2 の本当の快・不快を process-2 に入れたときの他者側の特徴予測の誤差で調べた．
誤差は経験なし条件 0.25〜0.29，経験50\%条件 0.24〜0.28 に対し，経験80\%条件では 0.38〜0.42 と大きかった（段階IIを行った組1〜4，二つのデータ）．
経験なし条件と経験50\%条件の値はほぼ同じであり，経験による転用が予測を正確にしたことは示されなかった．
経験80\%条件で零を入れるより悪かったのは，緑好みデータで4組中2組，好み混合データで3組であった．
この悪化は予想していなかった結果である．

A-2 の規則の確率を 0.5 から 0.8 にすると，経験80\%条件では $\Delta$ が小さくなり（本当の快・不快がより役立ち），その変化は経験50\%条件・経験なし条件より大きかった（図\ref{fig:cross}）．
組をまたいだ変化は経験80\%条件で $-0.096$ [$-0.139$, $-0.054$]（緑好み），$-0.098$ [$-0.125$, $-0.071$]（好み混合），経験50\%条件との交互作用は $-0.056$ [$-0.100$, $-0.012$]，$-0.046$ [$-0.072$, $-0.020$] で，事前登録した三つの予想は両データで6組すべて成り立った．
ただし，A-2 の確率を上げるとどの条件でも本当の快・不快はいくらか役立つようになり，悪化は消えずに小さくなった組が多く，組2 では零より悪いままであった．
また，確率を上げると A-2 も好きな色の角への滞在が増え（後半で好きな色から距離5未満にいる割合 0.99〜1.00），位置の分布も同時に変わる．
自分と他者の確率が一致したときに最もよく役立つことも示していない．
符号無作為条件でも，経験80\%条件の $\Delta$ そのものが組によってそろわず（悪化がはっきりしていたのは組2 だけ），悪化の原因は判定できなかった．
したがって，この結果は，自他の行動の仕方の違いが本当の快・不快の役立ち方に影響する可能性を示すにとどまり，自分の行動の仕方を他者に当てはめる投影と断定はできない．

\subsection{他者の好みの表現と推定}\label{sec:res-pref}
他者用モジュールは，構成を4通りに変えても（自由出力，形の制限，A-2 を隠した学習，1ビット）他者の好みを書き出さなかった（補足資料）．
事前登録した主の構成（経験50\%条件）では，T2・T5 とも4組中1組でしか合格の形にならず，その値（T2 $+0.03$〜$+0.05$，T5 $+0.04$〜$+0.05$）は本当の快・不快を入れた値（T2 $+0.60$，T5 $+0.42$）の 1/10 ほどであった．
一方，生成した $Q2$ を入れると他者側の特徴予測の誤差は本当の快・不快を入れた場合よりずっと小さく（0.116〜0.159 対 0.241〜0.276），モジュールは予測に役立つ位置の情報を伝えていた．
出力を1ビットに絞っても，好みと方角が食い違うエピソードで出力が好みに従う割合は 0.30〜0.41（経験80\%条件）で，方角の側に寄った．
これはモジュールの出力についての結果であり，内部状態に好みの情報がないことを示すものではない．

好みの候補を入れ比べる照合は，経験80\%条件の本番の3組で偶然を上回り（自己視界基準の BA 0.553〜0.701，事前登録した判定で合格），経験50\%条件では上回らなかった．
しかし方角だけの基準の BA は 0.708 [0.653, 0.759] で，自己視界基準ではモデルとの差が6組とも負（$-0.089$ [$-0.140$, $-0.038$]），他者特徴基準でも組をまたいでわずかに上回るにとどまった（$+0.049$ [$+0.026$, $+0.072$]，組ごとにはっきり上回ったのは組4 だけ）．
好みと方角が食い違うエピソードで選択が好みに従う基準を満たしたのは，他者特徴基準で6組中4組，自己視界基準で2組で，両方を満たしたのは組2 だけであった．
共有された予測の仕組みは，与えた候補をある程度区別できるが，他者固有の好みを独立に推定できることまでは示されなかった．

% ============================================================
% v2/sections_v2/discussion.tex  考察（短縮・再構成版，査読の指摘を受けて改訂 2026-10-27 17:00）
% 数字は sections_v2/results.tex と同じ出典。新しい数字は sections/discussion.tex（照合済み）から。
% ============================================================
\section{考察}\label{sec:discussion}

\subsection{何が新しく，何が分かったか}\label{sec:disc-what}
文献\cite{noguchi2022superposition}は，自他で重みを共有する予測学習から他者の位置・視点・運動の表現が生まれ，他者の観察から自分の運動に当たる信号が作られることを示した．
自分側で学んだ感覚運動の予測が他者側に使われうること自体は，そこですでに示されている．
本研究で新しく分かったのは，次の三点である．

第一に，与えた価値の信号に従って動いた経験から，報酬を用いない予測学習で価値に従う行動の見込みが生まれ，その見込みが，段階Iでは零しか入らない他者側に与えた快・不快にも反映された．

第二に，この転用は重みの共有だけでは生じなかった．
位置・運動・視界の系列を同じにして価値の信号と行動の対応だけを壊すと，自分側の快・不快への依存と他者側の符号に従う分は0前後になり，G5b も大きく弱まった．
この対照がなければ，経験80\%条件の効果は，赤の角に多くいた経験や，位置の手がかりとしての快・不快の大きさでも説明できた．
識別できたのは価値の信号と行動の一貫した対応の寄与であり，価値に従って動いた経験のあらゆる効果ではない．
G5b が経験なし条件より大きく残ったことは，赤への癖などの残りを含む．

第三に，学習後に写す比較から，転用は快・不快の入口の重みだけでは説明できず，自分の経験で学習された LSTM のほかの重みや計算にも依存する可能性が示された．
ただし，回復したのは視覚予測器を通らない量だけで，どの重みや計算が運ばれるかは特定していない．

予測に役立つ入力が使われるようになること自体は，予測学習の一般的な性質である．
本研究の知見は，その使い方が自分の行動の経験から学ばれて他者側の予測にも反映されること，そしてそれに要る条件の一つ（価値の信号と行動の一貫した対応）を，位置の分布をそろえた対照で識別したことにある．
これは，自分の意思決定の仕組みを他者の行動の仮説を試すために走らせるというシミュレーション説の考え\cite{gordon2021simulation}について，何がそろえば価値に関わる他者予測への転用が起きるかを，一つの構成の中で具体化したものと位置づけられる．
他の構成でも同じ条件が要るかは示していない．

\subsection{主張の強さと，前提を満たさない条件への依存}\label{sec:disc-strength}
経験80\%条件のモデルは，自己位置の追跡の事前基準を6組とも満たさなかった（共通テストで 0.70〜0.75，自分の学習データと同じ作りのテストでは 0.99）．
主張の強さは，これに応じて次のように分けられる．

前提を満たした経験50\%条件で支持されるのは，価値に従う行動の見込みが経験によって生まれること（6組すべてで経験なし条件との差）と，他者側の符号応答 G5b が事前登録した判定に合格したこと（本番の3組中2組）である．
ただし経験50\%条件の G5b は組をまたぐと差が0を含み（表\ref{tab:summary}），転用は弱く，組によってそろわない．

はっきりした転用，成立条件の切り分け（符号無作為条件，写す比較），他者の規則の確率を変えた比較は，すべて経験80\%条件に依っている．
このうち符号無作為条件は，経験80\%条件と同じ軌跡で学習し，自己位置の追跡も経験80\%条件に近かった．
したがって，両者の違いは前提の不成立では説明できない．
一方，経験80\%条件のモデルは評価に用いたランダムな歩きの場面に十分に一般化していないので，よく一般化したモデルで同じ大きさの効果が得られるかは分からない．
成立条件についての結論は，「赤の角にとどまる経験を多く積んだモデルで，価値の信号と行動の対応が転用に要った」という範囲で読むべきである．

\subsection{転用は他者の理解を保証しない}\label{sec:disc-st}
G5b と符号に従う分が示すのは，自分側で学んだ傾向が他者側の予測に反映されることであり，他者の行動の予測が正確になることではない．
実際，本当の快・不快を入れたときの誤差は経験なし条件と経験50\%条件でほぼ同じで，経験80\%条件ではかえって大きくなる組があった（\ref{sec:res-projection}節）．
自分由来の予測の転用は，他者の適切な理解を必ずしも保証しない．

人の他者理解では，自分を出発点にし，そこから他者に合わせて調整する過程が想定されている\cite{epley2004perspective,tamir2010neural}．
Suzuki ら\cite{suzuki2012learning}の区別でいえば，本研究のモデルは自分の価値づけを流用する経路を持つが，他者の行動の予測の外れから他者の価値を補正する経路を持たない．
他者の規則の確率を自分の学習条件に近づけると本当の快・不快がより役立ったことは，調整を欠いた転用が行動の仕方の違う他者で失敗する例と読むこともできる．
この読みで当てはめられているのは「何を好むか」ではなく「どれだけ強く好みに従うか」である．
しかし，規則の確率を変えると位置の分布も変わり，確率が一致したときに最もよく役立つことも示しておらず，符号無作為条件でも悪化の原因は判定できなかった．
この解釈は可能性にとどまる．
なお，自分の行為の経験が他者の行為の理解を変えること\cite{sommerville2005action}や，幼い子どもが他者に自分の好みを当てはめること\cite{repacholi1997early}と形の似た振る舞いも見られたが，これは計算モデルの振る舞いとの類推であり，人間について何かを示すものではない．

\subsection{なぜ他者の好みは取り出されないか}\label{sec:disc-shortcut}
他者用モジュールは他者側の予測の外れを減らすように学習するので，補正の学習信号そのものは与えられていた．
実際，経験80\%条件の上で学習したモジュールの出力を入れると，他者側の特徴予測の誤差は 0.112〜0.188 で，本当の快・不快を入れた場合（0.376〜0.418）よりずっと小さかった（組1〜4，二つのデータ）．
補正は学ばれたが，他者の価値を直す形ではなく，予測に役立つ位置や方角を伝える形をとった．
これは，学習の目的（予測の外れを減らすこと）はよく満たすが，意図した表現（他者の価値）とは別の手がかりに頼る近道\cite{geirhos2020shortcut}に当たると考えられる（推測であり，直接操作して確かめてはいない）．
本研究の課題では A-2 は一歩ごとに約半分しか規則に従わないので，本当の快・不快は A-2 の動きを部分的にしか決めず，位置や方角を伝える方が予測をずっとよくする．
なお，出力に好みが表れないことは，内部状態に好みの情報がないことを意味しない．

\subsection{限界}\label{sec:limitations}
第一に，環境は小さい．四隅にランドマークを置いた一つのアリーナで，カメラは回転しないので，「赤の方へ動く」ことは「赤の象限の方へ動く」ことと区別しにくい．
第二に，自他の対応は基準エンコーダで与え，価値は人が方角ごとに計算して与えた信号である．
本研究が示したのは，与えた価値の信号とそれに従う行動の規則性を予測学習が利用することまでである．
第三に，前提を満たさない条件への依存は\ref{sec:disc-strength}節で述べたとおりである．
第四に，非共有条件で他者側の働きが消えることは構成からほぼ予想され，写す比較は事前登録外で，予測画像を通る量では確かめられなかった．
第五に，符号無作為条件は行動の履歴そのものを変えていない．快・不快を常に零にした条件の結果は【結果待ち】である．
第六に，独立に学習した組は本番の3組，追加を含めて6組（符号無作為条件と写す比較は3組）と少なく，組ごとのばらつきも大きい．最終5点は互いに相関するチェックポイントで独立した繰り返しではなく，多重比較の補正も行っていない．符号に従う分（組1〜4）は事後に定義した指標である．
第七に，照合の候補は研究者が用意した二択であり，ネットワークが自ら他者の好みを推し量ったとは言えない．

\subsection{今後の課題}\label{sec:future}
\ref{sec:disc-shortcut}節の解釈からは，他者の行動の仕方が自分と同じなら，他者用モジュールは価値を書き出すようになるはずだという予測が出る．
本当の快・不快が役立つようになるという前半は\ref{sec:res-projection}節の比較で支持されたので，後半を確かめるには，A-2 が 0.8 の確率で規則に従うデータで他者用モジュールを学習する必要がある．
また，好きな物に着くとそれが価値を失い，次に好きな物へ向かう環境のように，好みの順位でしか次の行き先が決まらない課題では，位置や方角の近道が使えなくなるので，照合や他者用モジュールが方角を超えて好みを取り出すかを調べられる．
ただし，観察から順位を推す材料が本当にあるかを，学習の前にデータだけで確かめる必要がある．
そのほか，符号無作為条件の組を増やし，自己位置の追跡の前提を満たす条件でも同じ対照を行うこと，写した状態に合わせて視覚予測器を学習し直して予測画像を通る量でも確かめること，他者の行動の予測の外れから他者の価値を補正する経路を加えること，自他の対応と価値の表現を同時に学ぶエンコーダで同じ結果が得られるかを確かめることが挙げられる．

% ============================================================
% v2/sections_v2/conclusion.tex  おわりに（査読の指摘を受けて改訂 2026-10-27 17:00）
% ============================================================
\section{おわりに}\label{sec:conclusion}
本研究は，重ね合わせ機構\cite{noguchi2022superposition}の上で，自分の価値の経験から学んだ予測傾向がどのような条件で他者側の予測に転用されるかを調べた．
価値の信号に従って動いた経験からは，報酬も目標の入力も用いない予測学習によって価値に従う行動の見込みが生まれ，それは共有された重みを通して他者側に与えた快・不快にも反映された．
同じ軌跡で価値の信号と行動の対応だけを壊すとこの転用は大きく弱まり，他者の状態を追う力だけでは生じず，快・不快の入口の重みを写すだけでは戻らなかった．
一方，転用は他者の行動の予測を一般に正確にはせず，他者固有の好みの自発的な表現や推定も保証しなかった．
自分の経験に基づく予測は他者予測の出発点になりうるが，それだけでは他者の理解に足りない．
これらの条件の切り分けは，自己位置の追跡の前提を満たさなかった経験80\%条件のモデルで得たものであり，前提を満たすモデルでの確認が残る．


\section*{補足資料}
組ごとの値，学習後に重みを写す比較の詳細，真値条件の悪化の詳細，他者用モジュールの4通りの構成，仮説の照合の詳細，好み切替条件，基準エンコーダの選定，事前登録の経緯，関門の定義は補足資料に示す．

\section*{データとコードの公開}
実験に用いたシミュレータ，学習と評価のコード，設定ファイル（乱数のシードを含む），データの生成手順，図表を作るスクリプト，事前登録の記録は，\todo{【要確認：公開先と公開の時期（先生と相談）】}で公開する．

\bibliographystyle{unsrt}
\bibliography{references}

\end{document}
